\documentclass[11pt]{article}

% ---------------------------------------------------------------------------
% morisien-embed: text embedding models and evaluation for Mauritian Creole.
% Self-contained preprint source (compiles with a standard TeX distribution:
% pdflatex x2). Suitable for arXiv and Zenodo.
%
% Every number in this file is produced by a committed script and recorded in
% results/. The file that backs each claim is named in the text or in a table
% note, so a reader can check any figure without rerunning the pipeline.
% ---------------------------------------------------------------------------

\usepackage[a4paper,margin=1in]{geometry}
\usepackage[T1]{fontenc}
\usepackage[utf8]{inputenc}
\usepackage{lmodern}
\usepackage{microtype}
\usepackage{booktabs}
\usepackage{amsmath}
\usepackage{amssymb}
\usepackage{url}
\usepackage[round]{natbib}
\usepackage[hidelinks]{hyperref}
\usepackage{titlesec}

\titleformat{\section}{\normalfont\large\bfseries}{\thesection}{0.6em}{}
\titleformat{\subsection}{\normalfont\normalsize\bfseries}{\thesubsection}{0.6em}{}
\setlength{\parskip}{0.4em}
\setlength{\parindent}{0pt}

% Model names such as morisien-embed-v1.5 are long and admit few break points, which pushes a few
% lines past the margin. A small emergency stretch lets TeX loosen those lines instead.
\emergencystretch=1em

\title{\textbf{morisien-embed: Text Embedding Models and Evaluation\\for Mauritian Creole (Kreol Morisien)}}
\author{Singaraj B\\
\normalsize Independent Researcher\\
\normalsize \texttt{singaraj.yali@gmail.com}}
\date{2026}

\begin{document}
\maketitle

\begin{abstract}
Mauritian Creole (Kreol Morisien, ISO~639-3 \texttt{mfe}) is spoken at home by roughly 90\% of the
population of Mauritius, yet no dedicated text embedding model existed for it, and no available
evaluation could separate the multilingual models that cover it. We present \textbf{morisien-embed},
a sentence embedding model for Mauritian Creole, and \textbf{MorisienMTBitextMining}, the first
Mauritian Creole task in the Massive Text Embedding Benchmark (MTEB). On held-out MorisienMT bitext
retrieval the model reaches 0.9655 nDCG@10 against 0.9393 for LaBSE, the strongest general model we
evaluated, and a mean F$_1$ of 0.925 against 0.848 on the MTEB task. That comparison is in-domain.
To test whether it transfers, we build a harder pool from FLORES+ using the released xSIM++
augmentation: 996 queries over 45{,}029 passages, where each distractor is a gold passage with one
meaning-critical token changed. On that pool the margin disappears. morisien-embed scores 0.8439
nDCG@10 and untrained LaBSE 0.8449, a difference two paired tests cannot separate from zero. A
controlled ablation locates the result: hard-negative mining is worth $+0.0358$ out of domain while
Matryoshka loss contributes nothing to retrieval quality ($-0.0013$), and the choice of base model
matters more than any training decision. Fine-tuning LaBSE instead of multilingual-e5-base is worth
$+0.0171$ out of domain (McNemar $p = 0.0263$) and is the only configuration we measured that beats
untrained LaBSE; we release it as \textbf{morisien-embed-v1.5}. Both models fail in the same place:
distractors that reverse a causal relation are 4.4\% of the pool and cause about 51\% of the errors.
Models, benchmarks, and the full pipeline are released openly.
\end{abstract}

\section{Introduction}

Mauritian Creole is a French-lexified creole spoken at home by roughly 90\% of the population of
Mauritius (2022 census). Despite this, it is under-served by language technology: it is absent from
the widely used multilingual embedding evaluations, and general-purpose multilingual encoders,
trained with little or no Mauritian Creole supervision, embed it inconsistently. Text embedding
models underpin semantic search, retrieval-augmented generation, bitext mining, and clustering; the
lack of a strong Mauritian Creole embedder is a concrete barrier to building such systems for the
language.

The first version of this work reported a dedicated model that outperformed every general
multilingual encoder we could evaluate on Mauritian Creole bitext retrieval. That result was measured
on a single test set drawn from the corpus the model was trained on, because it was the only
Mauritian Creole retrieval evaluation available with any discrimination left in it. This version
reports what happened when we built one that had more.

This paper makes four contributions.
\begin{enumerate}
\itemsep0.1em
\item \textbf{An evaluation instrument.} A Mauritian Creole retrieval pool with measurable headroom,
built from FLORES+ using the released xSIM++ augmentation~\citep{chen2023xsimpp}, on which a purely
lexical baseline scores 0.3273 accuracy@1 and LaBSE 0.6928, against acceptance thresholds recorded
with the benchmark.
\item \textbf{Two models.} morisien-embed, a 278M-parameter model fine-tuned from
multilingual-e5-base, and morisien-embed-v1.5, a 471M-parameter model fine-tuned from LaBSE. The
first wins in domain and does not transfer; the second holds the in-domain result and is the only
configuration we measured that beats untrained LaBSE out of domain.
\item \textbf{A benchmark task.} MorisienMTBitextMining, the first Mauritian Creole task contributed
to MTEB, enabling standardized evaluation of any embedding model on the language, together with a
precise account of what it does and does not measure.
\item \textbf{A controlled ablation.} A measurement of what each element of the recipe is worth on
each benchmark, including the base model, corpus composition, and low-rank adaptation.
\end{enumerate}

We report the negative results in full, including a second benchmark we built that failed its own
acceptance criteria (Section~\ref{sec:smol}) and a hypothesis about the training corpus that the
measurement contradicted (Section~\ref{sec:composition}).

\section{Related Work}

Sentence embedding models map text to fixed-length vectors whose geometry reflects meaning; a common
recipe fine-tunes a pretrained encoder with a contrastive objective over paraphrase or translation
pairs~\citep{reimers2019sbert}. For multilingual and cross-lingual use, LaBSE~\citep{feng2022labse}
aligns 109 languages into a shared space and is a strong bitext-mining baseline, and the
multilingual-e5 family~\citep{wang2024me5} provides competitive general-purpose multilingual
encoders. Matryoshka representation learning~\citep{kusupati2022matryoshka} trains a single model
whose leading sub-vectors are themselves usable embeddings, allowing truncation for faster search at
a controlled accuracy cost. Low-rank adaptation~\citep{hu2022lora} trains a small number of injected
parameters instead of the whole encoder; \citet{biderman2024lora} characterize the trade as learning
less and forgetting less, which is the behaviour we observe in Section~\ref{sec:lora}. We evaluate
within the Massive Text Embedding Benchmark (MTEB)~\citep{muennighoff2023mteb}, the standard
framework for comparing embedding models across tasks and languages.

Evaluating bitext quality at scale motivated xSIM++~\citep{chen2023xsimpp}, which extends an
evaluation set with rule-based, hard-to-distinguish English perturbations and reports error rates by
perturbation type rather than a single aggregate. We use its released FLORES augmentation as the
basis of our out-of-domain pool. FLORES+ itself is the actively maintained successor to the FLORES
benchmark of \citet{nllb2024}, curated by the Open Language Data Initiative, and carries Mauritian
Creole.

Data for Mauritian Creole comes principally from MorisienMT~\citep{dabre2022morisienmt}, a
Creole$\leftrightarrow$\{English, French\} machine-translation dataset, and from
Kreyòl-MT~\citep{robinson2024kreyolmt}, a broad collection of parallel text for Latin American,
Caribbean, and colonial African creoles. SMOL~\citep{caswell2025smol} contributes professionally
translated data for 115 under-represented languages including Mauritian Creole; we hold it out of
training and report what happens when it is added (Section~\ref{sec:composition}). None of these
works targets embeddings; we repurpose their parallel text as contrastive supervision and reserve the
MorisienMT test split for evaluation.

\section{Data}
\label{sec:data}

We merge MorisienMT (MIT license) with the Mauritian Creole portion of Kreyòl-MT, giving 69{,}525 raw
rows and 35{,}064 unique Creole$\leftrightarrow$\{English, French\} pairs after exact deduplication
(22{,}211 English, 12{,}853 French). Kreyòl-MT carries mixed source licenses, so it is used for
training only and is not redistributed; we release the trained weights rather than the merged corpus,
and provide code to reconstruct it from the original sources at pinned revisions.

\textbf{The corpus is mostly a dictionary.} 22{,}164 of the 35{,}064 pairs, 63.21\%, have a
single-word Creole side, and the median Creole side is one word. Only 9{,}183 pairs reach six words.
This matters because every benchmark used here is built from sentences, so the training distribution
and the evaluation distribution differ sharply. Section~\ref{sec:composition} measures the
consequences, which are not the ones we expected.

\textbf{Public Creole text we do not train on.} The corpus is not all publicly available Mauritian
Creole parallel text. SMOL~\citep{caswell2025smol} carries 2{,}472 further gold pairs for the
language under CC-BY-4.0, with a median Creole length of 14 words. We hold them out of training
throughout and report an experiment that adds them in Section~\ref{sec:composition}.

\textbf{Leakage control.} Before training, every Creole sentence in the MorisienMT development and
test splits is looked for in the training pool under two keys: exact string match, and a
punctuation-, case-, and accent-insensitive match. Both are checks rather than removal steps. Running
them removes 0 of the 69{,}525 raw rows, because the MorisienMT evaluation splits and the training
sources are already disjoint; deduplication accounts for all 34{,}461 dropped rows. Stating that the
evaluation sentences are verifiably absent is a stronger claim than stating that a filter removed
them, and it is the one the measurement supports. Two residual effects remain and are discussed in
Section~\ref{sec:limitations}: target-side (English/French) texts are not reserved, and an audit
found one benchmark passage of 999 also present in training as the translation of a different Creole
sentence (its removal changes nDCG@10 by less than 0.0001). Corpus figures in this section are
produced by \texttt{scripts/audit\_corpus.py} and recorded in
\texttt{results/phase-d-corpus-audit.json}.

\section{Model and Training}
\label{sec:training}

morisien-embed is initialized from multilingual-e5-base (an XLM-RoBERTa-based encoder, 278M
parameters, 768-dimensional output, 512-token maximum) and fine-tuned in two stages. morisien-embed-v1.5
is the same recipe initialized from LaBSE (471M parameters, 768-dimensional output, 256-token
maximum).

\textbf{Hard-negative mining.} A stage-1 model is fine-tuned on all 35{,}064 pairs and used to mine
hard negatives with positive-aware false-negative filtering (5 negatives per anchor, a minimum rank
of 10, and a relative margin of 0.05). The margin filter is strict: roughly 24{,}100 of the 35{,}064
pairs retain a full negative set and form the contrastive training data. Every arm reported in this
paper mines from the same stage-1 checkpoint, so the seed variance we report covers training
initialization and not mining.

\textbf{Contrastive training.} The released checkpoints are trained on the mined tuples with a cached
multiple-negatives ranking loss (batch size 128, giving 767 in-batch negatives per anchor) wrapped in
a Matryoshka loss over dimensions $\{768, 512, 256, 128, 64\}$. Training runs for 3 epochs at
learning rate $2\times10^{-5}$ with 10\% warmup and fp16, seed 42, on a single T4 GPU (roughly 30
minutes of contrastive training plus 11 minutes of mining). No query or passage prompt is used at
training or inference. Dataset loads are pinned to exact revisions and model loads accept an explicit
revision, so the pipeline is reproducible: building the FLORES benchmark twice gives identical
checksums for queries, corpus, and relevance judgements.

\section{Benchmarks}

\subsection{In-domain: MorisienMT and the MTEB task}
\label{sec:mteb}

To make evaluation on Mauritian Creole standard and reproducible, we contribute
\textbf{MorisienMTBitextMining} to MTEB, the first Mauritian Creole task in the benchmark. It is a
bitext-mining task over the held-out MorisienMT test split, scored by F$_1$, hosted under the MTEB
organization at a pinned revision and available in the \texttt{mteb} package.

The task presents four directional subsets (mfe$\rightarrow$eng, eng$\rightarrow$mfe,
mfe$\rightarrow$fra, fra$\rightarrow$mfe) of 1{,}000 pairs each. These are not four independent
samples. All four are built over the same 1{,}000 MorisienMT test rows, comprising 999 distinct
Creole sentences, with the English and French sides supplying the other half of each pair. A model's
four scores are therefore four views of one set of Creole sentences, and they should be read as such.

Two properties of that set bear on how hard it is. First, its sentences are long: every one has at
least 10 words, with a median of 15 and a maximum of 49, while the training corpus has a median of
one word and only 14.08\% of its rows reach 10 words. Second, its wording is largely familiar. Under
punctuation-insensitive tokenization, 97.01\% of the Creole test tokens already appear in the
training corpus, and 74.10\% of the test sentences contain no token the model has not seen in
training. (Counting raw whitespace tokens instead, the figures are 92.83\% and 45.60\%.) The task
measures retrieval over held-out sentences, not over held-out vocabulary.

\subsection{Out of domain: a pool with headroom}
\label{sec:xsim}

The first version of this work evaluated out-of-domain generalization on the FLORES+ \texttt{mfe}
devtest split, 1{,}012 professionally translated sentences with no overlap with the training corpus.
Both morisien-embed (1.0000 accuracy@1) and LaBSE (0.9996) sat at the ceiling, and we read this as
evidence of no out-of-domain degradation rather than as a margin. That reading was too generous to
the pool. A character n-gram TF-IDF baseline with no neural model scores 0.8221 accuracy@1 on it, so
retrieving the right passage among 1{,}012 needs little more than surface overlap. The saturation was
a property of the pool, not of the language.

We therefore built a harder pool. xSIM++~\citep{chen2023xsimpp} releases rule-based English-side
augmentations of FLORES in which each distractor is a gold sentence with one semantically critical
element changed. We align the released \texttt{eng\_Latn} devtest augmentation to FLORES+ by exact
string match (996 of 996 originals align) and use the Creole side as queries. The resulting pool has
996 queries over 45{,}029 passages, of which 44{,}033 are distractors under three rules:
entity mention replacement (38{,}855, 88.24\%), number replacement (3{,}262, 7.41\%), and causality
alternation (1{,}916, 4.35\%).

Two acceptance conditions, recorded with the benchmark, decide whether it is usable: a lexical
baseline had to score below 0.50 accuracy@1, and a strong general model had to sit below the ceiling.
The same TF-IDF control scores 0.3273 and LaBSE 0.6928, so both hold. Construction, the exact control specification,
and checksums are recorded in \texttt{results/phase-a-xsim-eng.json}.

\subsection{A benchmark that failed its own criteria}
\label{sec:smol}

We also built a second held-out pool from SMOL, drawing distractors from the FLORES passage set. It
passes the lexical condition (TF-IDF 0.3643) and fails the second: LaBSE scores 0.9740 accuracy@1,
leaving no headroom. The reason is a construction error. Because the distractors come from a
different corpus than the gold passages, a model can eliminate them on domain alone without reading
the sentence. xSIM++ avoids this because every distractor is itself a near-copy of a gold passage. We
report this because only the first condition had been checked before the pool was built and
evaluated, and because a pool constructed this way will flatter any model put in front of it. It is
recorded in \texttt{results/phase-a2-smol.json} and is not used for any claim in this paper.

\section{Results}

\subsection{In domain}
\label{sec:indomain}

Table~\ref{tab:retrieval} reports Creole$\rightarrow$English retrieval on the held-out MorisienMT
test split. morisien-embed reaches 0.9655 nDCG@10 against 0.9393 for LaBSE, the strongest baseline,
and far ahead of larger general encoders such as multilingual-e5-large. morisien-embed-v1.5 holds
that result to within 0.0001.

\begin{table}[h]
\centering
\small
\begin{tabular}{lrrr}
\toprule
Model & Params & nDCG@10 & acc@1 \\
\midrule
paraphrase-multilingual-MiniLM-L12-v2 & 118M & 0.16 & 0.10 \\
BAAI/bge-m3 & 568M & 0.46 & 0.36 \\
intfloat/multilingual-e5-small & 118M & 0.54 & 0.42 \\
intfloat/multilingual-e5-base & 278M & 0.6402 & 0.53 \\
intfloat/multilingual-e5-large & 560M & 0.73 & 0.65 \\
sentence-transformers/LaBSE & 471M & 0.9393 & 0.91 \\
\textbf{morisien-embed} & \textbf{278M} & \textbf{0.9655} & \textbf{0.9440} \\
\textbf{morisien-embed-v1.5} & \textbf{471M} & \textbf{0.9654} & \textbf{0.9460} \\
\bottomrule
\end{tabular}
\caption{Creole$\rightarrow$English retrieval on the MorisienMT test split (1{,}000 queries, 999
passages). Figures are quoted at the precision at which they were recorded; the four models evaluated
only for the first version are recorded to two decimals.}
\label{tab:retrieval}
\end{table}

Repeating the contrastive stage with three seeds over the same deterministically mined negatives
gives 0.9661 nDCG@10 (standard deviation 0.0013) for the morisien-embed recipe and 0.9655 (0.0005)
for morisien-embed-v1.5. Both released checkpoints are seed 42, fixed in advance. This spread
measures sensitivity to initialization with the mined negatives held fixed; it is not the uncertainty
on the score, which we quantify in Section~\ref{sec:basemodel}. morisien-embed also leads in the
other two directions: Creole$\rightarrow$French, 0.9751 against LaBSE's 0.9475;
English$\rightarrow$Creole (999 queries over 1{,}000 passages), 0.9588 against 0.9247. Matryoshka
truncation degrades gracefully: nDCG@10 is 0.9591 at 256 dimensions and 0.9531 at 128. None of these
three has been rerun for morisien-embed-v1.5.

Table~\ref{tab:mteb} reports F$_1$ on the four directional subsets of the MTEB task, subject to the
qualification in Section~\ref{sec:mteb} that all four are views of one set of 999 Creole sentences.
morisien-embed attains a mean F$_1$ of 0.925 and leads every subset; LaBSE averages 0.848.

\begin{table}[h]
\centering
\small
\begin{tabular}{lrrrrr}
\toprule
Model & mfe$\to$eng & eng$\to$mfe & mfe$\to$fra & fra$\to$mfe & mean \\
\midrule
intfloat/multilingual-e5-small & 0.358 & 0.454 & 0.475 & 0.495 & 0.446 \\
sentence-transformers/LaBSE & 0.882 & 0.845 & 0.886 & 0.779 & 0.848 \\
\textbf{morisien-embed} & \textbf{0.927} & \textbf{0.909} & \textbf{0.939} & \textbf{0.924} & \textbf{0.925} \\
\bottomrule
\end{tabular}
\caption{F$_1$ on MTEB MorisienMTBitextMining. morisien-embed-v1.5 has not been evaluated on this
task.}
\label{tab:mteb}
\end{table}

\subsection{Out of domain}
\label{sec:ood}

On the xSIM++ pool the in-domain ordering does not hold. Table~\ref{tab:ood} gives the comparison.
morisien-embed scores 0.8439 nDCG@10 and untrained LaBSE 0.8449. Two paired tests on the same query
set agree that this difference is not measurable: McNemar's exact test on accuracy@1 gives
$p = 0.6852$ (106 wins for the model, 113 for LaBSE), and a paired bootstrap over queries gives a
difference of $-0.0010$ with a 95\% interval of $[-0.0160, +0.0135]$. The margin demonstrated in
domain does not transfer.

\begin{table}[h]
\centering
\small
\begin{tabular}{lrr}
\toprule
Model & nDCG@10 & acc@1 \\
\midrule
character n-gram TF-IDF (no neural model) & n/a & 0.3273 \\
intfloat/multilingual-e5-base & 0.6691 & 0.4639 \\
sentence-transformers/LaBSE & 0.8449 & 0.6928 \\
morisien-embed & 0.8439 & 0.6857 \\
\textbf{morisien-embed-v1.5} & \textbf{0.8616} & \textbf{0.7239} \\
\bottomrule
\end{tabular}
\caption{Creole$\rightarrow$English retrieval on the xSIM++ pool (996 queries, 45{,}029 passages).
The lexical control is an acceptance check on the pool, not a retrieval system, so only its
accuracy@1 is reported.}
\label{tab:ood}
\end{table}

\section{Ablations}

All ablations are evaluated on both benchmarks. Every arm uses the recipe of
Section~\ref{sec:training} with one element changed, mines from the same stage-1 checkpoint, and runs on
one T4 GPU. Results are in \texttt{results/phase-b-ablation.json} and
\texttt{results/phase-c-release-decision.json}.

\subsection{What each element of the recipe contributes}

Table~\ref{tab:ablation} builds the recipe up one element at a time. The two benchmarks disagree
about where the value is, which is the point of running both.

\begin{table}[h]
\centering
\small
\begin{tabular}{llrrr}
\toprule
Arm & Change & Seeds & In-domain & Out-of-domain \\
\midrule
R0 & stage 1, batch 48, no mining & 1 & 0.9535 & 0.8028 \\
R1 & + batch 128, cached loss & 1 & 0.9589 & 0.8081 \\
R2 & + hard-negative mining & 1 & 0.9663 & 0.8439 \\
R3 & + Matryoshka (morisien-embed) & 3 & 0.9661 & 0.8426 \\
R4 & LaBSE base (morisien-embed-v1.5) & 3 & 0.9655 & 0.8597 \\
R5 & LoRA $r{=}16$ on multilingual-e5-base & 1 & 0.9115 & 0.7957 \\
R6 & LoRA $r{=}16$ on LaBSE & 3 & 0.9553 & 0.8628 \\
\bottomrule
\end{tabular}
\caption{Component ablation, nDCG@10 on both benchmarks. Multi-seed arms report the mean; standard
deviations out of domain are 0.0032 (R3), 0.0021 (R4), and 0.0001 (R6).}
\label{tab:ablation}
\end{table}

Out of domain, hard-negative mining with the cached loss carries the result at $+0.0358$. Raising the
batch from 48 to 128 is worth $+0.0053$. Matryoshka loss is worth $-0.0013$, that is, nothing for
retrieval quality, though it still buys the truncation reported in Section~\ref{sec:indomain}. In domain
the same three are $+0.0074$, $+0.0054$, and $-0.0002$. The first version credited three techniques
for its result; one of them carries it.

\subsection{The base model}
\label{sec:basemodel}

Changing the initialization from multilingual-e5-base to LaBSE, holding everything else fixed, is
worth $+0.0171$ nDCG@10 out of domain and $-0.0006$ in domain. This is the largest single effect we
measured, larger than any training choice. It is supported by two paired tests on seed 42 of each
arm: McNemar's exact test gives $p = 0.0263$ (135 wins to 100), and a paired bootstrap over queries
gives $+0.0154$ with a 95\% interval of $[+0.0011, +0.0300]$. It is also supported by the seed
evidence, where the lowest LaBSE-based run (0.8574) is above the highest e5-based run (0.8461).

Two cautions apply to the interval. Run-to-run drift on an identical command reaches 0.0024 on this
pool, which is larger than the interval's lower bound, so $+0.0011$ should not be read as a floor on
the effect. The nondeterminism enters with hard-negative mining rather than with training: the R0 arm
reproduces exactly on a repeat, while R3 seed 42 gave 0.8444 and then 0.8461, and R4 seed 42 gave
0.8592 and then 0.8616. The seed spread, where the two arms' ranges do not overlap, is the stronger
evidence.

The size of the gain is itself informative. Measured against each arm's own untrained base,
fine-tuning is worth $+0.1735$ out of domain on multilingual-e5-base and $+0.0148$ on LaBSE. The
larger gain still lands below untrained LaBSE (0.8426 against 0.8449) while the smaller one lands
above it (0.8597). A large gain from a low floor is still a worse model. LaBSE's model card lists 109
languages; Mauritian Creole is not among them but Haitian Creole is, so its starting point on this
language is transfer from a related French-lexified creole rather than Creole supervision. The
Haitian probe in Section~\ref{sec:limitations} measures the same relationship from the other
direction.

\subsection{Corpus composition and volume}
\label{sec:composition}

Section~\ref{sec:data} noted that 63.21\% of the training corpus is single Creole words while every
benchmark is sentences. The obvious inference is that the dictionary entries dilute the training
signal. The measurement contradicts it.

Removing them, training on only the 9{,}183 pairs of six words or more, scores 0.8547 out of domain.
A volume-matched control, 9{,}183 pairs drawn at random from the full corpus, scores 0.8613. Three
seeds each, with non-overlapping ranges. The dictionary entries help rather than dilute: removing
them costs 0.0066. The volume-matched control also scores 0.0024 above the full corpus, which equals
the measured run-to-run drift and should be read as no difference, so within this recipe composition
matters and volume does not.

Adding data in the other direction fails as well. Adding SMOL's 2{,}468 sentence-level pairs, which
lowers the single-word share from 63.21\% to 59.06\%, costs 0.0048 out of domain on one seed, placing
it below the worst of the three seeds it is compared against, and changes the in-domain score by
$+0.0006$. Both directions of changing the corpus make the model worse, which suggests the corpus
sits near a local optimum for this recipe rather than being straightforwardly improvable. These are
recorded in \texttt{results/phase-c2-dictionary-ablation.json} and
\texttt{results/phase-c5-smol-training.json}. Caveats are noted there: one draw of the random subset,
one seed for the SMOL arm, and SMOL grew the mining pool by 7\%, so the negatives shifted along with
the data.

\subsection{Low-rank adaptation}
\label{sec:lora}

An $r{=}16$ LoRA adapter on LaBSE trains 3.27M parameters, 0.7\% of the total once the adapter is
attached, and gives the highest out-of-domain score we measured, 0.8628, with a seed standard
deviation of 0.0001 against 0.0021 for the same recipe fine-tuned in full. It is the most reproducible arm in the study. It costs 0.0102 in domain
against the full fine-tune, which is the trade \citet{biderman2024lora} describe as learning less and
forgetting less. We report it as an ablation rather than releasing it: its 0.0031 out-of-domain gain
over the full fine-tune is close to the 0.0024 run-to-run drift, and it pays for that with an
in-domain loss an order of magnitude larger. The same adapter on multilingual-e5-base is the worst
arm in the study at 0.7957, so the benefit belongs to the combination of method and base model rather
than to the method.

\section{Error Analysis}

Reporting the xSIM++ error rate, the fraction of queries whose nearest passage is not the gold one,
makes these numbers comparable to other xSIM++ results. morisien-embed scores 0.3143, LaBSE 0.3072,
and multilingual-e5-base 0.5361, agreeing with the nDCG comparison in Section~\ref{sec:ood} that the
first two are close.

The breakdown by perturbation rule is the finding, and it is the same for both strong models.
Causality alternation is 4.35\% of the distractor pool and causes 50.8\% of morisien-embed's errors
and 51.3\% of LaBSE's, an enrichment of 11.7 and 11.8 times its share. Entity mention replacement is
88.24\% of the pool and causes 29.4\% and 32.7\% of the errors respectively, an enrichment of about
0.3. Number replacement sits between them at roughly 2.2 to 2.7 times. Both models handle a
substituted entity and fail on a reversed causal direction, and they fail on it to the same degree.
A dedicated Creole fine-tune does not change which distinctions the encoder can make.

One caveat qualifies the diagnostic. Reading the distractor text shows that entity replacement often
produces a sentence no fluent reader would accept, for example passage \texttt{d3380}:
``Abderrahman Bouanane Flow is the study of the movement of individual drivers and vehicles between
two points and the interactions they make with one another.'' A model can reject such a passage
without representing its meaning, which is the most likely reason this rule is so under-represented
among the errors. The causality result does not depend on that reading, but the entity result should
not be taken as evidence that the models understand entities. Enrichment here is computed against
each rule's share of all 44{,}033 distractors rather than per query, so it measures which rule the
errors land on, not a per-rule error probability. Full figures are in
\texttt{results/phase-a-xsim-error-rate.json}.

\section{Limitations}
\label{sec:limitations}

\textbf{The margin over LaBSE is either in-domain or small.} morisien-embed's advantage over LaBSE
holds on the MorisienMT test split, which is drawn from the corpus it was trained on, and is not
measurable on the xSIM++ pool. morisien-embed-v1.5 does beat untrained LaBSE out of domain, by 0.0148
using the three-seed mean and 0.0167 using the released checkpoint, against a run-to-run drift of
0.0024. That margin is consistent across three seeds with no overlap. It is not large.

\textbf{Recipe selection had test visibility.} During recipe development the held-out test score was
printed at the end of every training run, so the choice of recipe was made with that score visible.
The first version of this paper bounded the resulting optimism at approximately 0.01 nDCG. We
withdraw that bound: no committed script produces it and the procedure was not specified well enough
to repeat. The exposure is stated plainly instead. The three-seed replications and every ablation in
this paper were run after the recipe was frozen. FLORES+ ships a 997-sentence \texttt{dev} split for
Mauritian Creole that no code in this repository has ever loaded; it remains an untouched holdout for
future recipe selection.

\textbf{Two uncertainties, not one.} The seed spread we report (0.0005 to 0.0032 depending on arm)
measures sensitivity to initialization with mined negatives held fixed. The sampling uncertainty on a
score, from the paired bootstrap over queries, is roughly $\pm 0.010$, and run-to-run drift on an
identical command reaches 0.0024. Differences smaller than these should not be read as effects. The
first version reported a seed interval of $\pm 0.0002$ without this distinction, which understated
the uncertainty on its headline number.

\textbf{What the in-domain benchmark measures.} As set out in Section~\ref{sec:mteb}, the four MTEB
subsets are views of one set of 999 Creole sentences, those sentences are long relative to the
training corpus, and 97.01\% of their tokens appear in training. Bitext retrieval over 999 passages
is also an easier problem than retrieval over a web-scale collection, so absolute scores here are
upper bounds on what the same model would achieve at scale.

\textbf{Haitian Creole proximity.} Like every multilingual embedder we tested, the model places
Haitian Creole close to Mauritian Creole. On 1{,}012 aligned mfe/hat/eng FLORES+ triplets, injecting
every same-meaning Haitian twin into the corpus drops mfe$\rightarrow$eng accuracy@1 from 1.00 to
0.68 (LaBSE resists this better, 0.79); asked instead to distinguish the two creoles, the fine-tune
picks Mauritian 709 times out of 1{,}012 against LaBSE's 351. Wrong-meaning Haitian text is never
confused, but mixed mfe/hat corpora will degrade retrieval.

\textbf{Not measured for morisien-embed-v1.5.} Creole$\rightarrow$French retrieval, the
English$\rightarrow$Creole direction, the MTEB task, and Matryoshka truncation quality are reported
for morisien-embed only. Its maximum sequence length is 256 tokens, LaBSE's default, against 512 for
morisien-embed.

\textbf{Other limitations.} Accuracy@1 near 0.944 in domain means roughly one query in eighteen ranks
a wrong translation first. Fine-tuning multilingual-e5-base costs some pure-English quality (STS-b
Spearman $\approx$ 0.79 against the base model's $\approx$ 0.85), so a general model is preferable for
English-only work. ALL-CAPS text embeds measurably differently from its lower-case form (cosine
$\approx$ 0.81). The available data over-represents religious, political, and literary registers, and
mixes pre- and post-2011 orthographies.

\section{Availability}

The models (\texttt{Singaraj/morisien-embed} and \texttt{Singaraj/morisien-embed-v1.5}) and the
benchmark dataset are released on the Hugging Face Hub; the MorisienMTBitextMining task is available
in the \texttt{mteb} package; and the complete data-construction, training, and evaluation code, from
which every number in this paper is reproducible, is at
\url{https://github.com/LK-maker-007/morisien-embed}. Each result file under \texttt{results/}
records the commit, environment, and command that produced it. The models and the benchmark dataset
are released under the MIT license.

\section{Conclusion}

We set out to build a dedicated embedding model for Mauritian Creole and to show it beat general
multilingual encoders. In domain it does, by a clear margin, and the MTEB task makes that comparison
standard and repeatable for any model. Out of domain, on a pool built to have headroom, it does not:
morisien-embed and untrained LaBSE are statistically indistinguishable. What the ablation shows
is that the choice of base model dominates every training decision we made, that hard-negative mining
carries what gain there is, and that the corpus sits near a local optimum that neither pruning nor
augmentation improved. Fine-tuning LaBSE instead, released as morisien-embed-v1.5, is the only
configuration we measured that beats untrained LaBSE out of domain, and the margin is small.

The more useful outcome may be the instrument rather than the model. Mauritian Creole now has a
retrieval evaluation on which models differ, and it shows both the fine-tune and its baseline failing
on the same distinction, reversed causal direction, at nearly twelve times its share of the pool.
That is a target for the next model rather than a limit of this one.

\begin{thebibliography}{9}
\small

\bibitem[Biderman et al.(2024)]{biderman2024lora}
Dan Biderman, Jacob Portes, Jose Javier Gonzalez Ortiz, Mansheej Paul, Philip Greengard, Connor
Jennings, Daniel King, Sam Havens, Vitaliy Chiley, Jonathan Frankle, Cody Blakeney, and John P.
Cunningham.
\newblock LoRA Learns Less and Forgets Less.
\newblock \emph{Transactions on Machine Learning Research}, 2024. arXiv:2405.09673.

\bibitem[Caswell et al.(2025)]{caswell2025smol}
Isaac Caswell, Elizabeth Nielsen, Jiaming Luo, Colin Cherry, Geza Kovacs, Hadar Shemtov, Partha
Talukdar, Dinesh Tewari, Baba Mamadi Diane, Djibrila Diane, Solo Farabado Cissé, Koulako Moussa
Doumbouya, Edoardo Ferrante, Alessandro Guasoni, Christopher Homan, Mamadou K. Keita, Sudhamoy
DebBarma, Ali Kuzhuget, David Anugraha, Muhammad Ravi Shulthan Habibi, Sina Ahmadi, Anthony Munthali,
Jonathan Mingfei Liu, and Jonathan Eng.
\newblock SMOL: Professionally Translated Parallel Data for 115 Under-represented Languages.
\newblock In \emph{Proceedings of the Tenth Conference on Machine Translation (WMT)}, pages
1103--1123, 2025. arXiv:2502.12301.

\bibitem[Chen et al.(2023)]{chen2023xsimpp}
Mingda Chen, Kevin Heffernan, Onur Çelebi, Alex Mourachko, and Holger Schwenk.
\newblock xSIM++: An Improved Proxy to Bitext Mining Performance for Low-Resource Languages.
\newblock In \emph{Proceedings of the 61st Annual Meeting of the Association for Computational
Linguistics (ACL)}, 2023. arXiv:2306.12907.

\bibitem[Dabre and Sukhoo(2022)]{dabre2022morisienmt}
Raj Dabre and Aneerav Sukhoo.
\newblock MorisienMT: A Dataset for Mauritian Creole Machine Translation.
\newblock \emph{arXiv preprint arXiv:2206.02421}, 2022.

\bibitem[Feng et al.(2022)]{feng2022labse}
Fangxiaoyu Feng, Yinfei Yang, Daniel Cer, Naveen Arivazhagan, and Wei Wang.
\newblock Language-agnostic BERT Sentence Embedding.
\newblock In \emph{Proceedings of the 60th Annual Meeting of the Association for Computational
Linguistics (ACL)}, 2022. arXiv:2007.01852.

\bibitem[Hu et al.(2022)]{hu2022lora}
Edward J. Hu, Yelong Shen, Phillip Wallis, Zeyuan Allen-Zhu, Yuanzhi Li, Shean Wang, Lu Wang, and
Weizhu Chen.
\newblock LoRA: Low-Rank Adaptation of Large Language Models.
\newblock In \emph{International Conference on Learning Representations (ICLR)}, 2022.
arXiv:2106.09685.

\bibitem[Kusupati et al.(2022)]{kusupati2022matryoshka}
Aditya Kusupati, Gantavya Bhatt, Aniket Rege, et al.
\newblock Matryoshka Representation Learning.
\newblock In \emph{Advances in Neural Information Processing Systems (NeurIPS)}, 2022.
arXiv:2205.13147.

\bibitem[Muennighoff et al.(2023)]{muennighoff2023mteb}
Niklas Muennighoff, Nouamane Tazi, Lo\"{i}c Magne, and Nils Reimers.
\newblock MTEB: Massive Text Embedding Benchmark.
\newblock In \emph{Proceedings of the 17th Conference of the European Chapter of the Association for
Computational Linguistics (EACL)}, 2023. arXiv:2210.07316.

\bibitem[NLLB Team(2024)]{nllb2024}
NLLB Team.
\newblock Scaling neural machine translation to 200 languages.
\newblock \emph{Nature}, 630(8018):841--846, 2024.

\bibitem[Reimers and Gurevych(2019)]{reimers2019sbert}
Nils Reimers and Iryna Gurevych.
\newblock Sentence-BERT: Sentence Embeddings using Siamese BERT-Networks.
\newblock In \emph{Proceedings of the 2019 Conference on Empirical Methods in Natural Language
Processing (EMNLP)}, 2019. arXiv:1908.10084.

\bibitem[Robinson et al.(2024)]{robinson2024kreyolmt}
Nathaniel R. Robinson et al.
\newblock Krey\`{o}l-MT: Building MT for Latin American, Caribbean and Colonial African Creole
Languages.
\newblock In \emph{Proceedings of the 2024 Conference of the North American Chapter of the
Association for Computational Linguistics (NAACL)}, 2024. arXiv:2405.05376.

\bibitem[Wang et al.(2024)]{wang2024me5}
Liang Wang, Nan Yang, Xiaolong Huang, Linjun Yang, Rangan Majumder, and Furu Wei.
\newblock Multilingual E5 Text Embeddings: A Technical Report.
\newblock \emph{arXiv preprint arXiv:2402.05672}, 2024.

\end{thebibliography}

\end{document}
